% Planned tables and figures for the ACL paper.
% This file is a checklist. Move finalized LaTeX tables/figure commands here
% once the experimental numbers are locked.

\section*{Main Tables}

% Table 1: Benchmark and temporal shift summary.
% Columns:
% Benchmark | Domain | Tasks | Span | Ordering | Metric | Primary shift
% Rows:
% PolyBench | prediction markets | 5,075 | 16 days | market close time |
% traded accuracy/CWR return | liquidity, uncertainty, market regime
% CTF-Dojo | CTF cybersecurity | 261 | 2011--2024 | year/event/name |
% pass rate | challenge format, file count, runtime dependency
% FutureX | future event prediction | 332 | 82 days | resolution date |
% accuracy | language, searchability, difficulty, platform data

% Table 2: Main results across variants.
% This is the paper's primary empirical table.
% Columns:
% Benchmark | H0 Baseline | H1 Linear | Early Freeze | H5 Navigation |
% H5 Multi | Best | Regression vs H0 | Cost
% Use benchmark-appropriate main metric:
% PolyBench: traded accuracy and CWR return in same cell.
% CTF/FutureX: accuracy.
% Add deltas in parentheses.

% Table 3: Task-level transfer/regression analysis.
% Columns:
% Benchmark | Variant | Improved over H0 | Regressed vs H0 | Net |
% Late-stream accuracy | Artifact size | Branch count
% This table directly answers reviewer concerns about overfitting and
% whether navigation merely increases capacity.

% Table 4: Artifact audit.
% Columns:
% Artifact path | Type | Stationary/non-stationary | Evidence | Expected branch
% Include representative artifacts from prompts, skills, tools, infra, memory.
% Keep this table in appendix if space is tight.

\section*{Main Figures}

% Figure 1: Method overview.
% Draw one diagram with:
% (a) linear A-Evolve chain: main evolves through all batches and every task
% loads all artifacts;
% (b) decoupled strategy tree: main trunk plus specialized branches;
% (c) inference path: navigator reads branch leaves and routes each task;
% (d) evolution path: stationary failures deepen main, regime failures branch.
% This should be the main figure in Section 3.

% Figure 2: Temporal shift and performance over time.
% Multi-panel plot:
% x-axis = batch/time.
% y-axis left = benchmark score; y-axis right or shaded bands = shift property.
% Panels for PolyBench, CTF-Dojo, FutureX.
% Overlay H0, H1, early-freeze, H5, H5_multi if available.
% The visual should make RQ1 obvious: full linear evolution does not simply
% improve with more cycles under temporal shift.

% Figure 3: Routing behavior.
% Stacked area chart or alluvial diagram:
% x-axis = cycle/time; y-axis = fraction of tasks routed to each branch.
% Annotate branch creation points and major shift regions.
% A small side panel should show branch utilization by task property.

% Figure 4: Multi-agent evolution decomposition.
% Optional if space allows. Show planner -> assignments -> per-branch evolver
% calls -> commits/rebases/registered branches. Use in method or appendix.

\section*{Analysis Plots}

% Plot A: Regression/improvement scatter.
% x-axis = improvements over baseline; y-axis = regressions vs baseline.
% Each point = variant/benchmark. Good systems move down/right.

% Plot B: Artifact growth vs performance.
% x-axis = prompt chars or total artifacts; y-axis = score.
% Show that bigger workspaces are not automatically better, supporting the
% need for selective navigation.

% Plot C: Navigation confidence/calibration.
% If navigator confidence is logged, bin by confidence and report routing
% success proxy: branch pass rate or branch-property match.

% Plot D: Branch staleness.
% x-axis = cycles since last routed; y-axis = branch selection/pass rate.
% Use to motivate stale-branch filtering or future work.
